\chapter*{General Introduction}
\addcontentsline{toc}{chapter}{General Introduction}
Manufacturing execution systems (MES) are computerized systems used in manufacturing to track and document the transformation of raw materials to finished goods. MES provides information that helps manufacturing decision-makers understand how current conditions on the plant floor can be optimized to improve production output\hyperref[ref:MESdefinition]{[1]}.

Modern small-to-medium manufacturing workshops often face challenges such as delayed communication between planners and operators, manual paper-based status updates, inventory inconsistencies, and limited traceability of work history. To address these issues this project was carried out through a collaboration between the Faculty of Science of Monastir (FSM), and Kabelbaum Test Systems (KTS), the project focused on a scoped MES-style solution that implements work-order lifecycle management, stage tracking, basic inventory control,  real-time messaging, and an alert/notification subsystem, along with analitical and reporting tools for production monitoring and performance evaluation.

This report presents the design and implementation of a production-tracking platform developed as part of an end of study project. 






\textbf{Report Structure}

\begin{itemize}
    \item \textbf{Chapter 1 -- General context and preliminary study: } We will start by introducing the host company, along with the overall context of the project and its main issue. Next, we will examine and evaluate the current situation in order to identify its limitations and propose an effective solution. Finally, we will outline the methodology that will be followed to carry out this project.
    \item \textbf{Chapter 2 -- Analysis and Specifications of requirements:} This chapter focuses on defining both functional and non-functional requirements, as well as the project’s architecture. It concludes with the presentation of the overall use case and class diagrams, which serve as a foundation for developing the complete project backlog. 
    \item \textbf{Chapter 3 and 4 -- Release 1 and 2:} The following two chapters form the core of our report. These two chapters will be dedicated to the development of the two releases of our platform
    
    \item \textbf{General Conclusion:} Finally, we conclude with a general summary that presents a recap of the work completed and outlines some future directions.
\end{itemize}
